%%%%%%%%%%%%%%%%%%%%%%%%%%%%%%%%%%%%%%%%%%%%%%%%%%%%%%%%%%%%%%%%%%%%%%%%%%%%%%%%%%
\begin{frame}[fragile]\frametitle{}
\begin{center}
{\Large Directories}
\end{center}
\end{frame}

%%%%%%%%%%%%%%%%%%%%%%%%%%%%%%%%%%%%%%%%%%%%%%%%%%%%%%%%%%%%%%%%%%%%%%%%%%%%%%%%%%
\begin{frame}[fragile]\frametitle{Current Directory \& Listing Files}
The \texttt{os} module is the traditional way to interact with the filesystem: find out where
you are, move around, and see what's there:
\begin{lstlisting}
import os

print('Current dir:', os.getcwd())

os.chdir('..')                  # move up one directory
print('Now in    :', os.getcwd())

print('Entries   :', os.listdir('.'))   # files and subdirectories, unsorted
\end{lstlisting}
\end{frame}

%%%%%%%%%%%%%%%%%%%%%%%%%%%%%%%%%%%%%%%%%%%%%%%%%%%%%%%%%%%%%%%%%%%%%%%%%%%%%%%%%%
\begin{frame}[fragile]\frametitle{Checking and Building Paths}
The \texttt{os.path} submodule answers questions about a path and builds new ones portably
(using the right slash for the current OS):
\begin{lstlisting}
import os

p = os.path.join('data', 'reports', 'summary.csv')
print(p)                              # data/reports/summary.csv

print(os.path.exists(p))              # True/False
print(os.path.isdir('data'))          # True/False
print(os.path.isfile(p))              # True/False

print(os.path.basename(p))            # summary.csv
print(os.path.dirname(p))             # data/reports
print(os.path.abspath(p))             # full path from filesystem root
\end{lstlisting}
\end{frame}

%%%%%%%%%%%%%%%%%%%%%%%%%%%%%%%%%%%%%%%%%%%%%%%%%%%%%%%%%%%%%%%%%%%%%%%%%%%%%%%%%%
\begin{frame}[fragile]\frametitle{Modern Python: \texttt{pathlib}}
Since Python 3.4, \texttt{pathlib.Path} wraps the same functionality as \texttt{os.path} in an
object-oriented API -- paths are objects with methods, and \texttt{/} joins them:
\begin{lstlisting}
from pathlib import Path

p = Path('data') / 'reports' / 'summary.csv'   # / joins path components
print(p)                              # data/reports/summary.csv

print(p.exists())
print(p.parent)                       # data/reports
print(p.name)                         # summary.csv
print(p.suffix)                       # .csv
print(p.with_suffix('.json'))         # data/reports/summary.json

for entry in Path('.').iterdir():
    print(entry)
\end{lstlisting}
\end{frame}

%%%%%%%%%%%%%%%%%%%%%%%%%%%%%%%%%%%%%%%%%%%%%%%%%%%%%%%%%%%%%%%%%%%%%%%%%%%%%%%%%%
\begin{frame}[fragile]\frametitle{Intuition: \texttt{os.path} vs.\ \texttt{pathlib}}
\begin{block}{Intuition}
\texttt{os.path} treats a path as a plain string that you slice and dice by passing it through
function after function. \texttt{pathlib} treats a path as an \emph{object} you can ask
questions of and chain operations on (\texttt{p.parent.parent / "x"}) -- similar to how
f-strings replaced \texttt{\%}-formatting: the same capability, but more readable, chainable
code.
\end{block}
\end{frame}

%%%%%%%%%%%%%%%%%%%%%%%%%%%%%%%%%%%%%%%%%%%%%%%%%%%%%%%%%%%%%%%%%%%%%%%%%%%%%%%%%%
\begin{frame}[fragile]\frametitle{Walking a Directory Tree}
\lstinline|os.walk()| is a generator that recursively visits every subdirectory, yielding
\texttt{(dirpath, dirnames, filenames)} for each one -- useful for processing an entire tree
without writing the recursion yourself:
\begin{lstlisting}
import os

for dirpath, dirnames, filenames in os.walk('.'):
    for filename in filenames:
        if filename.endswith('.py'):
            print(os.path.join(dirpath, filename))
\end{lstlisting}
The \texttt{pathlib} equivalent for a simple recursive glob:
\begin{lstlisting}
from pathlib import Path

for p in Path('.').rglob('*.py'):
    print(p)
\end{lstlisting}
\end{frame}

%%%%%%%%%%%%%%%%%%%%%%%%%%%%%%%%%%%%%%%%%%%%%%%%%%%%%%%%%%%%%%%%%%%%%%%%%%%%%%%%%%
\begin{frame}[fragile]\frametitle{Creating and Removing Directories}
\begin{lstlisting}
import os, shutil

os.mkdir('output')                    # fails if 'output' already exists
os.makedirs('output/2026/reports')    # creates all intermediate dirs; no-op if it exists

os.rmdir('output/2026/reports')       # fails unless the directory is empty
shutil.rmtree('output')               # removes a directory and everything inside it
\end{lstlisting}
\lstinline|shutil.rmtree| is not undoable -- there is no ``move to trash'' -- so double-check
the path before calling it in real code.
\end{frame}

%%%%%%%%%%%%%%%%%%%%%%%%%%%%%%%%%%%%%%%%%%%%%%%%%%%%%%%%%%%%%%%%%%%%%%%%%%%%%%%%%%
\begin{frame}[fragile]\frametitle{Exercise}
Write a function \lstinline|count_py_files(root)| that recursively counts how many \texttt{.py}
files exist under a given directory \texttt{root}.
\end{frame}

%%%%%%%%%%%%%%%%%%%%%%%%%%%%%%%%%%%%%%%%%%%%%%%%%%%%%%%%%%%%%%%%%%%%%%%%%%%%%%%%%%
\begin{frame}[fragile]\frametitle{Solution}
\begin{lstlisting}
from pathlib import Path

def count_py_files(root):
    return sum(1 for _ in Path(root).rglob('*.py'))

print(count_py_files('.'))
\end{lstlisting}
\end{frame}

%%%%%%%%%%%%%%%%%%%%%%%%%%%%%%%%%%%%%%%%%%%%%%%%%%%%%%%%%%%%%%%%%%%%%%%%%%%%%%%%%%
\begin{frame}[fragile]\frametitle{Quick Check: Directories}
\begin{block}{Think About It}
\lstinline|os.walk('.')| and \lstinline|Path('.').rglob('*.py')| both find every \texttt{.py}
file recursively. What's the practical difference in how they do it?
\end{block}
\pause
\begin{block}{Answer}
\lstinline|os.walk()| yields \texttt{(dirpath, dirnames, filenames)} level by level -- you
filter \texttt{filenames} yourself, and can even prune \texttt{dirnames} in place to skip whole
subtrees. \lstinline|Path.rglob()| directly yields matching \texttt{Path} objects for a single
glob pattern -- more concise for ``find all files matching X'', but you'd reach for
\lstinline|os.walk()| when you need finer control, like skipping specific directories entirely.
\end{block}
\end{frame}